\def\level{Première}  % Classe à droite
\def\course{NSI}
\def\chaptername{Représentation des entiers relatifs} % Titre au centre
\def\docname{RD03-S}        % Identifiant à gauche

\pagestyle{cours_FP}
\makeheader{} % Affiche le bandeau



\begin{easybox}[Codage binaire des entiers relatifs]{}
	\begin{minipage}{0.68\linewidth}
\begin{itemize}
	\item le bit de poids fort des nombres négatifs est 1
	\item le bit de poids fort des nombres positifs est 0
	\item il y a autant de positifs que de négatifs (0 est positif)
	\item $n$ bits permettent de coder $2^n$ valeurs différentes, soit $2^{n-1}$ valeurs positives et $2^{n-1}$ valeurs négatives.
	\item les valeurs codées avec $n$ bits sont comprises entre $-2^{n-1}$ et $2^{n-1}-1$
\end{itemize}
	\end{minipage} \hfill
	\begin{minipage}{0.30\linewidth}
		\includegraphics[width=0.79\linewidth]{codageZ.png}
	
	\end{minipage}

\end{easybox}

\begin{easybox}[Complément à 2~--~Première méthode]{}
	\begin{enumerate}
		\item Coder la valeur absolue du nombre en binaire
	   \item Inverser les bits de la représentation du nombre en base 2. (complément à 1)
	   \item Additionner 1 au résultat précédent. 
	\end{enumerate}
\end{easybox}



\begin{easybox}[Complément à 2~--~Deuxième méthode]{}
 \begin{enumerate}
 	\item Coder la valeur absolue du nombre en binaire
	\item Recopier les valeurs des bits de droite à gauche jusqu'au premier \verb|1|   inclus.
	\item Inverser les bits restant jusqu'au dernier (bit de gauche)
 \end{enumerate}
\end{easybox}

\bigskip

\def\level{Première}  % Classe à droite
\def\course{NSI}
\def\chaptername{Représentation des flottants} % Titre au centre
\def\docname{RD04-S}        % Identifiant à gauche

\makeheader{} % Affiche le bandeau

	\begin{methode}[Conversion d'un nombre décimal en base 2]{}
	\begin{enumerate}
	\item Conversion de la partie entière du nombre décimal
	\item Multiplication de la partie décimale par 2
	\item Récupération de la partie entière 
	\item Si la partie décimale est nulle on arrête sinon retour au 2.
	\item Assemblage de la partie entière et des bits récupérés pour former l'écriture binaire
\end{enumerate}
	\end{methode}

	\begin{methode}[Conversion d'un nombre binaire en décimal $100,0101_2$]{}

\begin{enumerate}
	\item Convertir la partie entière: $100_2=4_{10}$
	\item Convertir la partie décimale. On utilise un tableau des puissances de 2 négative.
\end{enumerate}

	
\renewcommand{\arraystretch}{1.3}
\begin{tabular}{|p{2.6cm}|p{1.3cm}|p{1.3cm}|p{1.3cm}|p{1.3cm}|p{1.3cm}|p{1.3cm}|}
\hline
Rang du bit & -1 & -2 & -3 & -4 & -5 & -6  \\ \hline
Puissance de 2 & $2^{-1}$ & $2^{-2}$ & $2^{-3}$ & $2^{-4}$ & $2^{-5}$ & $2^{-6}$  \\ \hline
Bit & 0 & 1 & 0 & 1 & 0 & 0  \\ \hline
Valeur & 0 & 0.25  & 0  & 0.0625 & 0 & 0 \\ \hline

\end{tabular}


\medskip

Il suffit ensuite d'additionner toutes les valeurs calculées: $100,0101_2=4+0.25+0.0625=4.3125_{10}$

\end{methode}
